
achève par
récurrence, en notant que si @@M0@@, et si @@M1@@ pour @@M2@@, on aura :
@@M3@@
d'où :
@@M4@@
d'où @@M5@@ d'après ce qui précède.

\section*{3. Conséquences du théorème de structure pour le @@M6@@ d'un fibré projectif}

\textbf{Proposition 3.1.} \uline{Soient} @@M7@@ \uline{un schéma quasi}-\uline{compact}, @@M8@@ \uline{des} @@M9@@\uline{-Modules localement libres de type fini}, @@M10@@ \uline{le schéma des drapeaux de} @@M11@@, @@M12@@.
\uline{Alors l'homomorphisme canonique}~:
@@M13@@
\uline{est injectif} (@@M14@@ \uline{étant défini}, \uline{soit à partir des} @@M15@@\uline{-Modules localement libres de type fini}, \uline{soit à partir des complexes parfaits}).

On vérifie immédiatement en utilisant la propriété universelle des schémas de drapeaux que @@M16@@ peut être construit comme suit : on considère la suite de @@M17@@-schémas @@M18@@, définis par récurrence en prenant pour @@M19@@ le schéma de drapeaux de l'image inverse de @@M20@@ sur @@M21@@; alors @@M22@@. On voit alors tout de suite que pour démontrer la proposition, il suffit de la montrer dans le cas où @@M23@@, le cas général s'en déduisant par récurrence. On peut donc supposer que @@M24@@ est le schéma des drapeaux d'un @@M25@@-Module localement libre de type fini @@M26@@.

On sait que @@M27@@ peut alors être construit comme suit : on pose @@M28@@, @@M29@@, et on définit par récurrence des schémas @@M30@@ et
